\documentclass[UTF8,a4paper,12pt]{ctexart}

\usepackage{geometry}
\geometry{left=2.6cm,right=2.6cm,top=2.5cm,bottom=2.5cm}

\usepackage{amsmath,amssymb}
\usepackage{graphicx}
\usepackage{float}
\usepackage{booktabs}
\usepackage{array}
\usepackage{fancyvrb}
\usepackage{caption}
\usepackage{hyperref}

\pagestyle{plain}
\pagenumbering{arabic}

\title{\textbf{脑与认知科学大作业报告}\\神经网络实现编解码}

\author{
班级：\underline{\makebox[4cm][c]{人工智能2402班}}
\quad
学号：\underline{\makebox[4.5cm][c]{202431121013030}}
\quad
姓名：\underline{\makebox[3cm][c]{胡弈周}}
}

\date{}
\begin{document}

\maketitle

\begin{abstract}
本文基于 Python 语言手动实现了一个 $8$-$3$-$8$ 结构的自编码器神经网络，用于完成 $8 \times 8$ 单位矩阵的编解码任务。网络由输入层、隐藏层和输出层组成，其中输入层和输出层均包含 $8$ 个神经元，隐藏层包含 $3$ 个神经元。程序没有使用 TensorFlow、PyTorch 等深度学习框架，而是使用 NumPy 完成前向传播、误差反向传播和梯度下降参数更新，并使用 Matplotlib 绘制训练过程中的损失函数曲线。实验结果表明，经过训练后，网络能够将 $8 \times 8$ 单位矩阵压缩到 $8 \times 3$ 的隐藏层编码表示，并较好地恢复出原始单位矩阵。
\end{abstract}

\newpage

\section{作业任务说明}

本次作业要求实现一个 $8$-$3$-$8$ 的自编码器神经网络结构。网络输入为 $8 \times 8$ 单位矩阵，每一行为一个训练样本，也就是一个 one-hot 向量。网络需要通过隐藏层完成编码，再通过输出层完成解码，使输出结果尽可能接近原始输入。

本实验程序使用 Python 编写，主要依赖 NumPy 和 Matplotlib。其中 NumPy 用于矩阵计算和手动实现梯度下降，Matplotlib 用于绘制损失函数图像。整个训练过程没有调用 TensorFlow、PyTorch 等深度学习框架，满足手动实现神经网络训练算法的要求。

\section{网络结构设计}

本实验采用三层全连接神经网络：

\begin{itemize}
    \item 输入层：$8$ 个神经元，对应 $8$ 维 one-hot 输入向量；
    \item 隐藏层：$3$ 个神经元，用于学习输入数据的压缩编码；
    \item 输出层：$8$ 个神经元，用于将隐藏层编码结果解码回 $8$ 维输出向量。
\end{itemize}

设输入矩阵为 $X \in \mathbb{R}^{8 \times 8}$，真实标签矩阵为 $Y \in \mathbb{R}^{8 \times 8}$。由于本任务是自编码任务，因此有：
\begin{equation}
    Y = X = I_8
\end{equation}

其中 $I_8$ 表示 $8 \times 8$ 单位矩阵。输入层到隐藏层的权重矩阵为 $W^{(1)} \in \mathbb{R}^{8 \times 3}$，偏置为 $b^{(1)} \in \mathbb{R}^{1 \times 3}$；隐藏层到输出层的权重矩阵为 $W^{(2)} \in \mathbb{R}^{3 \times 8}$，偏置为 $b^{(2)} \in \mathbb{R}^{1 \times 8}$。

\section{激活函数}

网络隐藏层和输出层均使用 Sigmoid 激活函数。Sigmoid 函数定义如下：
\begin{equation}
    \sigma(z)=\frac{1}{1+e^{-z}}
\end{equation}

其导数可以表示为：
\begin{equation}
    \sigma'(z)=\sigma(z)(1-\sigma(z))
\end{equation}

在程序实现中，为了防止指数计算时出现数值溢出，代码中使用了 \verb|np.clip(z, -60, 60)| 对输入值进行截断处理。

\section{前向传播过程}

前向传播过程分为输入层到隐藏层、隐藏层到输出层两个阶段。

首先，输入矩阵 $X$ 经过第一层线性变换，得到隐藏层的预激活值：
\begin{equation}
    Z^{(1)} = XW^{(1)} + b^{(1)}
\end{equation}

然后通过 Sigmoid 激活函数得到隐藏层输出：
\begin{equation}
    H = \sigma(Z^{(1)})
\end{equation}

其中 $H \in \mathbb{R}^{8 \times 3}$，即每个输入样本被编码为 $3$ 维隐藏层表示。

接着，隐藏层输出经过第二层线性变换，得到输出层预激活值：
\begin{equation}
    Z^{(2)} = HW^{(2)} + b^{(2)}
\end{equation}

最后通过 Sigmoid 激活函数得到网络预测输出：
\begin{equation}
    O = \sigma(Z^{(2)})
\end{equation}

其中 $O \in \mathbb{R}^{8 \times 8}$，表示网络解码后的输出矩阵。

\section{损失函数}

本实验采用均方误差作为损失函数。设训练样本数量为 $N$，输出维度为 $8$，则损失函数定义为：
\begin{equation}
    L = \frac{1}{2N}\sum_{n=1}^{N}\sum_{k=1}^{8}(O_{nk}-Y_{nk})^2
\end{equation}

程序中的实现为：
\begin{equation}
    L = 0.5 \times \operatorname{mean}\left(\sum (O-Y)^2\right)
\end{equation}

其中乘以 $0.5$ 是为了在反向传播求导时抵消平方项产生的系数 $2$，使梯度表达更加简洁。

\section{反向传播算法}

反向传播的目标是计算损失函数对各层权重和偏置的梯度，然后使用梯度下降法更新参数。

\subsection{输出层误差项}

输出层误差项为：
\begin{equation}
    \Delta^{(2)} = (O-Y) \odot O \odot (1-O)
\end{equation}

其中 $\odot$ 表示矩阵对应位置元素相乘。根据链式法则，输出层权重和偏置的梯度为：
\begin{equation}
    \frac{\partial L}{\partial W^{(2)}} =
    \frac{H^T\Delta^{(2)}}{N}
\end{equation}

\begin{equation}
    \frac{\partial L}{\partial b^{(2)}} =
    \operatorname{mean}(\Delta^{(2)})
\end{equation}

\subsection{隐藏层误差项}

隐藏层误差需要由输出层误差反向传播得到：
\begin{equation}
    \Delta^{(1)} =
    \left(\Delta^{(2)}(W^{(2)})^T\right)
    \odot H \odot (1-H)
\end{equation}

因此，输入层到隐藏层权重和偏置的梯度为：
\begin{equation}
    \frac{\partial L}{\partial W^{(1)}} =
    \frac{X^T\Delta^{(1)}}{N}
\end{equation}

\begin{equation}
    \frac{\partial L}{\partial b^{(1)}} =
    \operatorname{mean}(\Delta^{(1)})
\end{equation}

\section{梯度下降参数更新}

设学习率为 $\eta$，则参数更新公式如下：
\begin{equation}
    W^{(2)} \leftarrow W^{(2)} - \eta \frac{\partial L}{\partial W^{(2)}}
\end{equation}

\begin{equation}
    b^{(2)} \leftarrow b^{(2)} - \eta \frac{\partial L}{\partial b^{(2)}}
\end{equation}

\begin{equation}
    W^{(1)} \leftarrow W^{(1)} - \eta \frac{\partial L}{\partial W^{(1)}}
\end{equation}

\begin{equation}
    b^{(1)} \leftarrow b^{(1)} - \eta \frac{\partial L}{\partial b^{(1)}}
\end{equation}

本实验在每个 epoch 中对完整的 $8$ 个样本进行一次批量梯度下降更新。

\section{程序实现说明}

本实验代码主要包括四个逻辑部分：参数配置、数学工具函数、自编码器模型类和主程序训练流程。虽然最终提交代码可以保存在一个 \verb|brain_cognitive_autoencoder.py| 文件中，但从结构上可以理解为以下四个部分：

\begin{table}[H]
    \centering
    \caption{程序逻辑模块划分}
    \begin{tabular}{lll}
        \toprule
        模块 & 对应内容 & 主要功能 \\
        \midrule
        \verb|config.py| & 参数配置部分 & 设置输入维度、隐藏层维度、训练轮数、学习率等 \\
        \verb|math_utils.py| & 数学函数部分 & 实现 Sigmoid 函数和均方误差损失函数 \\
        \verb|autoencoder.py| & 模型定义部分 & 实现前向传播、反向传播、编码、解码和预测 \\
        \verb|main.py| & 主程序部分 & 构造数据、训练模型、保存图像和输出结果 \\
        \bottomrule
    \end{tabular}
\end{table}

其中，自编码器类 \verb|AutoEncoder838| 是程序的核心。其 \verb|forward| 函数用于完成前向传播，\verb|train_one_epoch| 函数用于完成一次完整的前向传播、损失计算、反向传播和参数更新。

\section{实验参数设置}

本实验采用的主要超参数如下表所示：

\begin{table}[H]
    \centering
    \caption{实验超参数设置}
    \begin{tabular}{ll}
        \toprule
        参数名称 & 参数取值 \\
        \midrule
        输入层神经元数 & $8$ \\
        隐藏层神经元数 & $3$ \\
        输出层神经元数 & $8$ \\
        训练轮数 & $50000$ \\
        学习率 & $10.0$ \\
        随机种子 & $42$ \\
        激活函数 & Sigmoid \\
        损失函数 & 均方误差 \\
        优化方法 & 批量梯度下降 \\
        \bottomrule
    \end{tabular}
\end{table}

权重初始化采用类似 Xavier 均匀初始化的方法。对于某一层输入维度为 $d_{in}$、输出维度为 $d_{out}$ 的权重矩阵，初始化范围为：
\begin{equation}
    \left[-\sqrt{\frac{6}{d_{in}+d_{out}}},\sqrt{\frac{6}{d_{in}+d_{out}}}\right]
\end{equation}

这种初始化方式有助于避免初始权重过大或过小，从而提升训练稳定性。

\section{实验结果与分析}

程序运行结束后，会输出最终损失值、原始单位矩阵、隐藏层编码矩阵、解码后的输出矩阵以及经过 $0.5$ 阈值二值化后的矩阵。同时，程序会保存训练过程中的损失函数图像。

在本次实验参数设置下，最终损失值约为：
\begin{equation}
    L \approx 7.8142 \times 10^{-5}
\end{equation}

隐藏层编码矩阵示例如下：
\begin{Verbatim}[fontsize=\small]
[[0.9934 0.0038 0.9310]
 [0.0461 0.0097 0.9917]
 [0.9791 0.9834 0.0044]
 [0.0031 0.0060 0.0675]
 [0.0032 0.9849 0.9752]
 [0.0123 0.9316 0.0063]
 [0.8483 0.0081 0.0040]
 [0.9955 0.9963 0.9962]]
\end{Verbatim}

解码后的输出矩阵示例如下：
\begin{Verbatim}[fontsize=\small]
[[0.9917 0.0043 0.0000 0.0000 0.0000 0.0000 0.0038 0.0076]
 [0.0057 0.9931 0.0000 0.0048 0.0065 0.0000 0.0000 0.0000]
 [0.0000 0.0000 0.9919 0.0000 0.0000 0.0041 0.0042 0.0073]
 [0.0000 0.0045 0.0000 0.9919 0.0000 0.0056 0.0062 0.0000]
 [0.0000 0.0040 0.0000 0.0000 0.9915 0.0038 0.0000 0.0076]
 [0.0000 0.0000 0.0059 0.0042 0.0066 0.9927 0.0000 0.0000]
 [0.0066 0.0000 0.0061 0.0041 0.0000 0.0000 0.9922 0.0000]
 [0.0040 0.0000 0.0040 0.0000 0.0036 0.0000 0.0000 0.9888]]
\end{Verbatim}

将输出矩阵按照阈值 $0.5$ 进行二值化后，得到：
\begin{Verbatim}[fontsize=\small]
[[1 0 0 0 0 0 0 0]
 [0 1 0 0 0 0 0 0]
 [0 0 1 0 0 0 0 0]
 [0 0 0 1 0 0 0 0]
 [0 0 0 0 1 0 0 0]
 [0 0 0 0 0 1 0 0]
 [0 0 0 0 0 0 1 0]
 [0 0 0 0 0 0 0 1]]
\end{Verbatim}

由结果可以看出，解码后的矩阵在主对角线位置接近 $1$，在非主对角线位置接近 $0$。经过阈值化处理后，输出矩阵能够恢复为 $8 \times 8$ 单位矩阵。这说明该 $8$-$3$-$8$ 自编码器已经基本学会了对单位矩阵的编码与解码映射。

\section{损失函数曲线}

训练过程中，程序将每一轮 epoch 的损失值保存下来，并绘制为损失函数曲线。随着训练轮数增加，损失函数整体呈下降趋势，说明网络参数通过反向传播和梯度下降不断被优化。

\begin{figure}[H]
    \centering
    % 将下面的 loss_curve.png 改成你实际生成的损失函数图片路径。
    % 例如：figures/loss_curve.png 或 ./loss_curve.png
    \IfFileExists{loss_curve.png}{
        \includegraphics[width=0.85\textwidth]{loss_curve.png}
    }{
        \fbox{\parbox[c][6cm][c]{0.85\textwidth}{\centering loss_curve.png}}
    }
    \caption{训练过程中的损失函数曲线}
    \label{fig:loss_curve}
\end{figure}

从图1可以观察训练过程中损失函数的变化情况。若曲线逐渐下降并趋于稳定，说明模型的训练过程是有效的。

\section{结论}

本文使用 Python 手动实现了一个 $8$-$3$-$8$ 结构的自编码器神经网络。实验中没有使用 TensorFlow、PyTorch 等深度学习框架，而是通过 NumPy 完成矩阵运算、前向传播、反向传播和梯度下降更新。模型以 $8 \times 8$ 单位矩阵作为训练数据，通过 $3$ 个隐藏层神经元完成压缩编码，并最终在输出层恢复原始输入。

实验结果表明，经过 $50000$ 轮训练后，网络损失降低到较小水平，解码输出矩阵与原始单位矩阵高度接近。经过阈值化处理后，输出结果可以恢复为标准单位矩阵，说明该神经网络成功实现了对 $8 \times 8$ 单位矩阵的编解码功能。

\end{document}
